% ============================================
% Quantitative Fellowship Program Cover Letter
% Federal Reserve System
% ============================================

\documentclass[11pt,letterpaper]{letter}

\usepackage{geometry}
\usepackage{microtype}
\usepackage[hidelinks]{hyperref}

\geometry{left=25mm,right=25mm,top=25mm,bottom=25mm}

\signature{Decory Edwards}
\address{
Department of Economics \\
Johns Hopkins University \\
Baltimore, MD 21218 \\
\texttt{dedwar65@jh.edu}
}

\begin{document}

\begin{letter}{
Quantitative Fellowship Program Hiring Committee \\
Federal Reserve System
}

\opening{Dear Members of the Hiring Committee,}

I am writing to apply for the Quantitative Fellowship Program at the Federal Reserve System. I am a Ph.D.\ candidate in Economics at Johns Hopkins University, advised by Professors Christopher D.~Carroll, Nicholas W.~Papageorge, and Stelios Fourakis, and I expect to complete my degree by May 2026. My training emphasizes macroeconomics, quantitative modeling, and applied data analysis, with a focus on understanding risk, heterogeneity, and policy transmission in complex economic systems.

My research agenda centers on the sources and consequences of heterogeneity in economic outcomes. In my job market paper, I estimate persistent heterogeneity in individual returns to wealth using micro data from the Survey of Consumer Finances and embed these estimates into a structural macroeconomic model. I show that accounting for heterogeneity in returns substantially improves the model’s ability to match observed wealth concentration and alters the predicted response of the economy to aggregate shocks. This work reflects the same core challenges faced in supervisory and regulatory settings, where persistent differences across institutions, balance sheets, and risk exposures play a central role in stress scenarios and policy evaluation.

In a second project, I use longitudinal data from the Health and Retirement Study to examine how beliefs and interpersonal trust affect portfolio choices and realized returns. I document a systematic, non-linear relationship between trust and returns to wealth, highlighting how behavioral and informational frictions can amplify or mitigate risk exposure over time. This project reinforces my interest in combining large administrative or survey datasets with rigorous statistical methods to identify persistent risk factors that matter for financial stability.

The Quantitative Fellowship Program is particularly appealing to me because it combines rigorous quantitative analysis with direct policy relevance and public service. The opportunity to rotate across Reserve Banks, contribute to stress test modeling, portfolio analysis, and model validation, and interact with senior supervisors and regulated institutions aligns closely with my training and professional goals. I have extensive experience programming in Python and Stata, managing large datasets, and developing structural and reduced-form models. I am comfortable communicating technical results to non-technical audiences and translating quantitative findings into clear economic narratives that support decision making.

I am motivated by the Federal Reserve’s mission to promote a safe, sound, and efficient financial system, and I would value the opportunity to contribute my quantitative skills and macroeconomic perspective in this setting. I am enthusiastic about the rotational structure of the QFP, the emphasis on professional development, and the possibility of continuing in a full-time quantitative role within the Federal Reserve System upon completion of the program.

Thank you very much for your time and consideration. I would welcome the opportunity to discuss my application further.

\closing{Sincerely,}

\end{letter}
\end{document}
